%% sections/03_mitigation.tex
%% Module 3: Risk Mitigation and Policy Implications

\subsection{Structural Risk Mitigation Frameworks}

The quantitative findings presented in Sections~\ref{sec:intro}--\ref{sec:impact}
imply several structural interventions that could reduce the probability of systemic
market disruption:

\subsubsection{Index Provider Concentration Guardrails}

S\&P Dow Jones Indices and Nasdaq currently apply ad-hoc concentration rebalancing
when any single constituent exceeds 24\% or the top-5 collectively exceed 50\% of index
weight \citep{sp500methodology,nasdaq100methodology}. Given that both thresholds were
exceeded by mid-2024, a more dynamic constraint---such as a real-time HHI-based
rebalancing trigger---would smooth concentration-driven amplification.

A proposed guardrail: trigger mandatory rebalancing whenever $\HHI > 800$ for any
major index. This would limit the concentration amplifier $\alpha_{\HHI}$ to
approximately $1.5\times$ in a worst-case scenario, reducing tail PFE by an
estimated 35--40\%.

\subsubsection{Compute Derivatives as Macro Hedging Instruments}

The most structurally sound mitigation mechanism is the development of standardised
\textit{compute derivatives}---forward and options contracts written on the price of
AI compute capacity (typically quoted in \$/H100-hour or \$/PFLOP). Such instruments
would allow:

\begin{itemize}
  \item Hyperscalers to hedge against compute price deflation (falling GPU rental prices
        that reduce the economic justification for data centre construction);
  \item Institutional investors to express directional views on AI infrastructure
        utilisation rates without direct equity exposure;
  \item Central banks to monitor aggregate compute capacity utilisation as a leading
        indicator of AI sector stress, analogous to the use of forward-looking
        credit spread indices in fixed-income surveillance.
\end{itemize}

The primary barrier to market development is the absence of a standardised reference
rate for spot compute pricing. CoreWeave, Lambda Labs, and Microsoft Azure all
publish different spot prices for equivalent GPU capacity, creating a fragmented
price discovery landscape. A LIBOR-equivalent rate for AI compute---a ``COMPOR''---would
be a prerequisite for a functioning derivatives market.

\subsubsection{Regulatory Intervention Thresholds}

Based on our modelling, the Systemic Short scenario becomes self-reinforcing above a
$-20\%$ peak-to-trough NDX drawdown: at that threshold, index-tracking margin calls
from leveraged ETF products would generate forced selling at an estimated \$8B/day,
which exceeds the estimated market-making buffer at typical institutional spread widths.

We recommend that SEC and CFTC establish pre-approved circuit-breaker coordination
protocols for AI-sector-specific drawdowns above $-15\%$ (Level~1), $-20\%$ (Level~2),
and $-30\%$ (Level~3) in the Nasdaq-100, distinct from the current market-wide
circuit breakers which are calibrated to S\&P~500 moves of $-7\%$, $-13\%$, and $-20\%$.

\subsection{Investor-Level Hedging Strategies}

For institutional investors with significant AI equity exposure, three practical hedging
overlays follow directly from the VaR/PFE framework:

\begin{enumerate}
  \item \textbf{Long-dated NDX put spreads:} Purchasing 18-month at-the-money NDX put
        options funded by selling 40\% out-of-the-money puts creates a cost-efficient
        protection structure that covers the Systemic Short scenario without full
        premium outlay.

  \item \textbf{VIX call spreads:} The 2022 drawdown demonstrated that implied volatility
        on the Nasdaq reached 35--45 at peak stress from a pre-shock level of 18--22.
        A VIX call spread (long 30, short 55) structured for the Q3--Q4 2026 window
        provides convex payoff precisely during the IPO supply shock window.

  \item \textbf{Sector rotation into compute-infrastructure beneficiaries:}
        Utilities and industrial REITs with data-centre power grid exposure
        (e.g., Vistra Energy, NRG) benefit from hyperscaler Capex regardless of
        whether AI applications reach commercial inflection, providing a hedged long.
\end{enumerate}
